\chapter{Theoretical Foundation}
\label{chap:theoretical_foundation}

\section{DNA Data Storage}
\label{sec:dna_data_storage}
% TODO: Explain density, durability, technical limitations, synthesis cost, and reading/sequencing challenges.

\section{Digital to DNA Encoding}
\label{sec:digital_to_dna_encoding}
% TODO: Present block segmentation, binary-to-base conversion, and the necessity of sequence constraints.

\section{Errors and Biological Constraints}
\label{sec:errors_and_biological_constraints}
% TODO: Detail substitutions, insertions, deletions, erasures, invalid reads, GC content, and homopolymers.
% Distinguish clearly between channel errors and biological sequence acceptability constraints.

\section{Erasure Codes and Fountain Codes}
\label{sec:erasure_codes_and_fountain_codes}
% TODO: Introduce the difference between fixed-rate and rateless codes, paving the way for LT codes.

\section{LT Codes}
\label{sec:lt_codes}
% TODO: Present droplet generation, bipartite graph representation, Ideal Soliton Distribution, and Robust Soliton Distribution.

\section{Robust Soliton Distribution}
\label{sec:robust_soliton_distribution}
% TODO: Formally explain rho(d), tau(d), mu(d), and discuss the practical role of parameters c and delta on low-degree droplets and ripple maintenance.

\section{Peeling Decoder and Ripple}
\label{sec:peeling_decoder_and_ripple}
% TODO: Describe the iterative belief-propagation algorithm over GF(2), identification of degree-1 droplets, removal of resolved variables, and the role of the ripple.

\section{DNA Fountain Architecture}
\label{sec:dna_fountain_architecture}
% TODO: Present how DNA Fountain combines LT codes with seed-based headers, biological screening, and mechanisms for discarding invalid reads.

\section{Evaluation Metrics}
\label{sec:evaluation_metrics}
% TODO: Formally define metrics used in experiments: success rate, overhead, screening rejection rate, XOR operations, ripple size, encoding time, and decoding time.
